% TOP_PROBLEM: {"id": 7000016, "title": "Geometry of Curves and Surfaces — Problem 3.3", "category_id": 6, "category": {"id": 6, "name": "geometry", "display_name": "Geometry", "description": "Euclidean and non-Euclidean geometry, geometric structures.", "slug": "geometry", "order_index": 6, "created_at": "2026-07-31T15:26:25.671Z"}, "status": "open"}
% FIRST_POSTED: null
% ENTRY_KIND: statement_only
% Imported from: batch18.tex; UnsolvedMath ID 7000016
\documentclass[11pt]{article}
\usepackage[margin=1in]{geometry}
\usepackage{amsmath,amssymb,mathtools}
\usepackage{fontspec,unicode-math}
\setmainfont{Latin Modern Roman}
\setsansfont{Latin Modern Sans}
\setmathfont{Latin Modern Math}
\usepackage{xurl,hyperref}
\hypersetup{colorlinks=true,urlcolor=blue,linkcolor=blue}
\newcommand{\sref}[2]{\hyperlink{source:#1:#2}{[S#2]}}
\newcommand{\cataloguescope}[1]{\paragraph{Recorded mathematical question.}#1}
\begin{document}
% UnsolvedMath ID 7000016; source catalogue code AMR-069-0016
\setcounter{section}{7000015}
\section{Geometry of Curves and Surfaces — Problem 3.3}\label{7000016}
{\small\sffamily UnsolvedMath ID 7000016 · Geometry}
\subsection{Definitions and mathematical statement}

\paragraph{Shared notation.}$\mathbb N=\{1,2,\ldots\}$, and $\mathbb Z,\mathbb Q,\mathbb R,\mathbb C$ denote the integers, rationals, reals and complex numbers. Any other conventions are specified locally.


A convex polyhedron $P$ here is the boundary of the convex hull of finitely many points of $\mathbb R^3$ which do not all lie in a plane. Its edge graph consists of its vertices and edges. A spanning edge tree $T$ contains all vertices and is connected and has no cycles. Cutting along $T$ produces a compact polyhedral disk $P_T$. Its planar development is a map $u_T:P_T\to\mathbb R^2$ that is locally an isometry and preserves lengths on each face. An unfolding is simple if $u_T$ is one-to-one; in that case the developed surface has no overlaps. A polyhedron is unfoldable if some spanning edge tree gives a simple unfolding. The source explains that allowing any collection of edge cuts with connected planar development yields the same question. No affine deformation of $P$ is part of this definition. The surface carries its intrinsic length metric. A pseudo-edge is a distance-minimizing geodesic segment joining two vertices of $P$. A pseudo-edge graph $G$ is a three-connected graph embedded in $P$, with exactly the vertices of $P$, made of pseudo-edges and with convex faces: every face interior angle is less than $\pi$. Three-connected means that deleting any two graph vertices leaves the graph connected. Cutting along a spanning tree $T$ of $G$ gives a surface with a planar development $u_T$. The pair $(P,G)$ is unfoldable if at least one such development is injective.
\paragraph{Statement recorded in the supplied source.}
Does there exist a convex polyhedron $P$ with a pseudo-edge graph $G$ such that $(P,G)$ is not unfoldable, so that every spanning-tree cut of $G$ has a noninjective planar development? The cited author paper answers this existence question affirmatively with a 340-vertex example. \sref{7000016}{2}
\subsection{Short English statement}
Can shortest geodesic edges form a graph on a convex polyhedron for which every spanning-tree unfolding overlaps? The source constructs such an example.
\subsection{Sources}
\begin{description}
\item[\hypertarget{source:7000016:1}{S1}] ULAM.AI and the UnsolvedMath Contributors, \emph{UnsolvedMath: A Curated Collection of Open Mathematics Problems}, record 7000016 (AMR-069-0016). CC BY 4.0. \url{https://huggingface.co/datasets/ulamai/UnsolvedMath}.
\item[\hypertarget{source:7000016:2}{S2}] Nicholas Barvinok and Mohammad Ghomi, \emph{Pseudo-edge unfoldings of convex polyhedra}, arXiv:1709.04944v4 (2019), Introduction, pp. 1--2, including Theorem 1.1; Section 3.1, pp. 4--5. \url{https://arxiv.org/pdf/1709.04944v4}.
\end{description}
\label{end:7000016}
\end{document}
